\section{Discussion and Limitations}
\label{sec:discussion}
\label{sec:limitations}

\subsection*{Optimizer inadequacy is not objective misalignment}

The nested test changes what a poor deeper run means: a terminal objective
above an embedded shallower point is certified optimizer inadequacy under that
run's budget. It does not certify global success when passed. Conversely,
lower mean energy can accompany lower feasible mass, and the O3--O0 contrast
remains positive after both arms pass the diagnostic under each tested
optimizer. Optimization quality and choice of loss must therefore be assessed
separately.

\subsection*{O2 identifies unused feasible-concentration capacity}

O2 supplies the missing counterfactual: what can the same ansatz and matched
optimizer do when feasible mass itself is requested? Its improvement over O0
prevents the O0 shortfall from being attributed wholly to ansatz capacity.
The held-out H2 result then gives O3 a sharper interpretation than a simple
O3--O0 win: CVaR approached this exact-feasibility capacity control within the
frozen margin. O2 requires exact membership aggregation and is neither a
shot-efficient nor deployment-ready objective.

\subsection*{Depth adds capacity and optimization difficulty}

The $p=4$ subset illustrates why depth should not be read from one endpoint.
It improves terminal feasible mass while producing more $p=3\to4$ than
$p=2\to3$ nested failures. The tested 120--480-call trajectories do not show
a monotone failure reduction, but they cannot establish behavior under much
larger budgets, broad multi-start policies, or other optimizers. Depth expands
the variational family and can simultaneously make its classical navigation
less reliable.

\subsection*{Endpoint robustness exceeds training robustness}

Sampling fixed exact-trained endpoints recovers the objective ordering
reliably as shots increase. Retraining with a freshly sampled loss is a harder
problem: both shot-trained point estimates preserve the O3--O0 direction, yet
both graph intervals cross zero and the 1,000-shot condition is less
consistent. The evidence supports finite-shot endpoint robustness and only a
qualified, mixed training result; it supplies no device-noise or hardware-
readiness evidence.

\subsection*{Scope}

The benchmark uses synthetic layered DAGs, one additive resource, exact CPU
statevectors, and small classically tractable tasks. The exact classical
solver confirms that the study is a mechanism probe, not a quantum--classical
performance contest. Its conclusions concern uniform-initialized, transverse-
$X$ Penalty-X QAOA and do not extend to all constrained QAOA. The $m=20$
reversal and administratively censored $m=22$ cells preclude a global scaling
claim; the latter records this study's preregistered computational budget, not
a fundamental limit at 22 qubits.

Structured or feasibility-preserving approaches change the information and
resource model by moving cost into state preparation, mixer construction, or
constraint-aware operators; they are therefore contextual comparators rather
than matched baselines for the present Penalty-X attribution study. In
particular, no XY, Dicke, Grover, or path mixer is assumed to solve general
RCSP feasibility. The membership-query discussion likewise does not
lower-bound the rich-energy CVaR experiment. Real-network validity, hardware
noise, compilation, and independent external reproduction remain outside the
tested scope, and no quantum-advantage claim is made.
